\documentclass[12pt,a4paper]{article}

% ------------------------------------------------------------
% Encoding, language, typography
% ------------------------------------------------------------
\usepackage[utf8]{inputenc}
\usepackage[T1]{fontenc}
\usepackage[english]{babel}
\usepackage{lmodern}
\usepackage{microtype}

% ------------------------------------------------------------
% Page layout
% ------------------------------------------------------------
\usepackage[a4paper,margin=2.5cm]{geometry}
\usepackage{setspace}
\onehalfspacing

% ------------------------------------------------------------
% Mathematics
% ------------------------------------------------------------
\usepackage{amsmath}
\usepackage{amssymb}
\usepackage{amsfonts}
\usepackage{mathtools}
\usepackage{bm}
\usepackage{bbm}

% ------------------------------------------------------------
% Tables, lists, links, colour
% ------------------------------------------------------------
\usepackage{booktabs}
\usepackage{enumitem}
\usepackage{xcolor}
\usepackage[hidelinks]{hyperref}

% colour for the shaded remark boxes
\usepackage{tcolorbox}
\tcbuselibrary{skins}
\definecolor{remarkbg}{RGB}{245,247,250}
\definecolor{remarkframe}{RGB}{180,195,215}

% ------------------------------------------------------------
% Citations
% ------------------------------------------------------------
\usepackage[round,authoryear]{natbib}

% ------------------------------------------------------------
% Theorem-style environments
% ------------------------------------------------------------
\usepackage{amsthm}
\newtheorem{assumption}{Assumption}
\newtheorem{definition}{Definition}
\newtheorem{proposition}{Proposition}
\newtheorem{lemma}{Lemma}
\newtheorem{corollary}{Corollary}
\theoremstyle{remark}
\newtheorem{remark}{Remark}

% ------------------------------------------------------------
% Useful math commands
% ------------------------------------------------------------
\newcommand{\E}{\mathbb{E}}
\newcommand{\En}{\mathbb{E}_N}
\newcommand{\R}{\mathbb{R}}
\newcommand{\ind}{\mathbbm{1}}
\newcommand{\muY}{\mu_Y}
\newcommand{\muD}{\mu_D}
\newcommand{\muYhat}{\hat{\mu}_Y}
\newcommand{\muDhat}{\hat{\mu}_D}
\newcommand{\ehat}{\hat{e}}
\newcommand{\YRF}{\widetilde{Y}^{\,\mathrm{RF}}}
\newcommand{\DFS}{\widetilde{D}^{\,\mathrm{FS}}}

% ------------------------------------------------------------
% Shaded intuition box
% ------------------------------------------------------------
\newenvironment{intuition}[1]{%
  \begin{tcolorbox}[
    enhanced,
    colback=remarkbg,
    colframe=remarkframe,
    arc=3pt,
    boxrule=0.6pt,
    left=8pt, right=8pt, top=6pt, bottom=6pt,
    title={\small\textit{Intuition: #1}},
    fonttitle=\bfseries\small,
    coltitle=black
  ]
  \small
}{%
  \end{tcolorbox}
}

% ------------------------------------------------------------
\title{%
  Double Robustness of the Wald-AIPW Estimator:\\
  Derivation, Proofs, and Interpretation%
}
\author{}
\date{}
% ------------------------------------------------------------

\begin{document}

\maketitle
\tableofcontents
\vspace{1em}

% ============================================================
\section{Setup and Target}
\label{sec:setup}
% ============================================================

Throughout, let $(Y_i, D_i, Z_i, X_i)_{i=1}^N$ be an i.i.d.\ sample where $Y_i$ is the
outcome, $D_i \in \{0,1\}$ the endogenous binary treatment, $Z_i \in \{0,1\}$ a binary
instrument, and $X_i$ a vector of pre-treatment covariates. The instrument is valid
conditional on $X_i$ in the sense of Assumptions~1--5 of the main text.

\paragraph{Target quantity.}
Define the instrument-conditional conditional expectations
\begin{align}
  \muY(z, x) &:= \E[Y_i \mid Z_i = z,\, X_i = x], \label{eq:muY} \\
  \muD(z, x) &:= \E[D_i \mid Z_i = z,\, X_i = x], \label{eq:muD} \\
  e(x)        &:= \Pr(Z_i = 1 \mid X_i = x). \label{eq:e}
\end{align}
The reduced-form and first-stage causal effects of the instrument are
\begin{align}
  \gamma_Y &:= \E[\muY(1, X_i) - \muY(0, X_i)], \label{eq:gammaY} \\
  \gamma_D &:= \E[\muD(1, X_i) - \muD(0, X_i)], \label{eq:gammaD}
\end{align}
and the LATE is the ratio $\tau_{\mathrm{LATE}} = \gamma_Y / \gamma_D$
\citep{imbens1994identification, tan2006distributional}.

\paragraph{The estimation problem.}
All three nuisance functions $\muY$, $\muD$, $e$ are unknown. Any estimator of
$\gamma_Y$ or $\gamma_D$ must therefore use estimated nuisance parameters
$\muYhat$, $\muDhat$, $\ehat$ in place of the truth. The central question is: what
happens to the estimate when these replacements are imperfect?

% ============================================================
\section{Three Strategies for Estimating the Reduced Form}
\label{sec:three_strategies}
% ============================================================

It is sufficient to focus on $\gamma_Y$; the treatment of $\gamma_D$ is identical
with $D_i$ replacing $Y_i$ throughout. There are three natural strategies, each
relying on a different nuisance parameter.

% ------------------------------------------------------------
\subsection{Strategy 1: Regression Adjustment (Wald-RA)}
\label{subsec:ra}
% ------------------------------------------------------------

\paragraph{Estimator.}
The regression-adjustment (RA) estimator predicts the potential outcome for each
individual under each instrument value using the outcome model, then averages the
contrast:
\begin{align}
  \hat{\gamma}_Y^{\,\mathrm{RA}}
  = \frac{1}{N} \sum_{i=1}^{N}
    \Bigl[\muYhat(1, X_i) - \muYhat(0, X_i)\Bigr].
  \label{eq:ra_estimator}
\end{align}

\paragraph{What it relies on.}
For each individual $i$, the estimator asks: given their covariate profile $X_i$, what
does the model predict their outcome would be under $Z=1$ versus $Z=0$? The prediction
$\muYhat(z, X_i)$ is a fitted value from a model trained on observations with $Z_i = z$.
The individual-level contrast $\muYhat(1,X_i) - \muYhat(0,X_i)$ is then a model-based
imputation of the counterfactual difference.

\paragraph{Identification: why it works when the model is correct.}
Under the conditional independence assumption $Z_i \perp\!\!\!\perp Y_i(z) \mid X_i$,
\begin{align}
  \E[\muY(1, X_i) - \muY(0, X_i)]
  = \E\bigl[\E[Y_i \mid Z_i=1, X_i] - \E[Y_i \mid Z_i=0, X_i]\bigr]
  = \gamma_Y.
  \label{eq:ra_id}
\end{align}
Hence $\hat{\gamma}_Y^{\,\mathrm{RA}} \xrightarrow{p} \gamma_Y$ whenever
$\muYhat(z,\cdot) \xrightarrow{p} \muY(z,\cdot)$ for both $z \in \{0,1\}$.

\paragraph{Failure mode: model misspecification.}
Suppose the true outcome surface is nonlinear but $\muYhat$ is a linear model. Let
$b(z, x) := \muYhat(z,x) - \muY(z,x)$ denote the pointwise bias of the model. Then
\begin{align}
  \hat{\gamma}_Y^{\,\mathrm{RA}}
  \;\xrightarrow{p}\;
  \E\bigl[\muY(1,X_i) - \muY(0,X_i)\bigr]
  + \underbrace{\E\bigl[b(1,X_i) - b(0,X_i)\bigr]}_{\text{model bias} \;\neq\; 0}.
  \label{eq:ra_bias}
\end{align}
The bias $\E[b(1,X_i) - b(0,X_i)]$ does \emph{not} shrink to zero with sample size
if the model is structurally misspecified. More data does not help: the estimator
consistently estimates the wrong quantity.

\begin{intuition}{Why a linear model for a curved surface goes wrong}
Imagine $Y_i$ and $X_i$ have a U-shaped relationship: outcomes are high for low-$X$
and high-$X$ individuals, but low in the middle. If the Z=1 group happens to contain
more high-$X$ individuals (because $Z$ is only conditionally random), a linear model
overpredicts $\muY(1,X)$ in the tails and underpredicts in the middle. The bias in the
Z=1 surface and the Z=0 surface do not cancel, so the contrast $\muYhat(1,X) -
\muYhat(0,X)$ is wrong everywhere. Averaging over $X_i$ does not fix this.
\end{intuition}

% ------------------------------------------------------------
\subsection{Strategy 2: Inverse Probability Weighting (Wald-IPW)}
\label{subsec:ipw}
% ------------------------------------------------------------

\paragraph{Estimator.}
Rather than modelling the outcome, the IPW estimator reweights the observed outcomes
by the inverse of the instrument propensity score:
\begin{align}
  \hat{\gamma}_Y^{\,\mathrm{IPW}}
  = \frac{1}{N} \sum_{i=1}^{N}
    \left[
      \frac{Z_i\, Y_i}{\ehat(X_i)}
      -
      \frac{(1-Z_i)\, Y_i}{1 - \ehat(X_i)}
    \right].
  \label{eq:ipw_estimator}
\end{align}

\paragraph{What it relies on.}
The weights $1/e(X_i)$ and $1/(1-e(X_i))$ rebalance the two instrument groups so
that, after reweighting, both groups have the same marginal distribution over $X_i$.
This removes the confounding from the covariate imbalance between $\{Z_i=1\}$ and
$\{Z_i=0\}$ without ever modelling the outcome surface.

\paragraph{Identification: why it works when the propensity score is correct.}
For the $Z=1$ term, conditioning on $X_i$ and using
$\E[\ind\{Z_i=1\} \mid X_i] = e(X_i)$:
\begin{align}
  \E\!\left[\frac{Z_i\, Y_i}{e(X_i)}\right]
  &= \E\!\left[\frac{\E[Z_i\, Y_i \mid X_i]}{e(X_i)}\right]
  = \E\!\left[\frac{e(X_i)\,\E[Y_i \mid Z_i=1, X_i]}{e(X_i)}\right]
  = \E[\muY(1, X_i)].
  \label{eq:ipw_id_z1}
\end{align}
The $e(X_i)$ cancels exactly. An identical argument gives
$\E[(1-Z_i)Y_i/(1-e(X_i))] = \E[\muY(0,X_i)]$, so
$\E[\hat{\gamma}_Y^{\,\mathrm{IPW}}] = \gamma_Y$ when $\ehat = e$.

\paragraph{Failure mode: propensity score misspecification.}
If $\ehat(x) \neq e(x)$, the cancellation in \eqref{eq:ipw_id_z1} is broken. Define
$\delta(x) := \ehat(x) - e(x)$ as the pointwise propensity score error. Then
\begin{align}
  \E\!\left[\frac{Z_i\, Y_i}{\ehat(X_i)}\right]
  = \E\!\left[\frac{e(X_i)\,\muY(1, X_i)}{\ehat(X_i)}\right]
  = \E\!\left[\muY(1,X_i)\cdot\frac{e(X_i)}{e(X_i)+\delta(X_i)}\right]
  \;\neq\; \E[\muY(1,X_i)]
  \label{eq:ipw_bias}
\end{align}
whenever $\delta \not\equiv 0$. The bias is the covariance between $\muY(1,X_i)$ and
the ratio $e(X_i)/\ehat(X_i)$: if the model underestimates $e(X_i)$ precisely for
high-$Y$ individuals, the upward bias in their weights produces an upward bias in the
estimate.

There is a second, more dangerous failure mode: \emph{weight instability}. When
$\ehat(X_i) \approx 0$ or $\ehat(X_i) \approx 1$, the weights
$1/\ehat(X_i)$ or $1/(1-\ehat(X_i))$ explode, amplifying small outcome values
to enormous scale. A single observation with a badly estimated propensity score can
dominate the entire sample average.

\begin{intuition}{Why wrong propensity scores are hard to detect}
Consider a near-randomised experiment where $e(X_i) = 0.03$ for a small group.
A logistic regression predicts $\ehat(X_i) = 0.05$ instead---a difference of only
0.02 in probability. But the IPW weight is $1/0.03 \approx 33$ with the true score
and $1/0.05 = 20$ with the estimated score---a 40\% difference in weight. Multiplied
by the outcome value, even a moderate outcome observation can shift the sample mean
substantially. The propensity score model may look excellent by any standard fit
statistic ($R^2$, Brier score) yet produce large weight distortions where it matters most.
\end{intuition}

% ============================================================
\section{The Augmented IPW: Combining Both Strategies}
\label{sec:aipw}
% ============================================================

% ------------------------------------------------------------
\subsection{Construction of the Pseudo-Outcome}
\label{subsec:pseudo}
% ------------------------------------------------------------

The AIPW estimator \citep{tan2006distributional} augments the outcome model prediction
with an IPW-weighted residual correction. Define the \emph{reduced-form pseudo-outcome}
for each individual as
\begin{align}
  \YRF_i
  &:=
  \underbrace{
    \muYhat(1, X_i) - \muYhat(0, X_i)
  }_{\text{(A) outcome model contrast}}
  \;+\;
  \underbrace{
    \frac{Z_i\bigl(Y_i - \muYhat(1, X_i)\bigr)}{\ehat(X_i)}
    -
    \frac{(1-Z_i)\bigl(Y_i - \muYhat(0, X_i)\bigr)}{1 - \ehat(X_i)}
  }_{\text{(B) IPW-weighted residual correction}}.
  \label{eq:rf_pseudo_main}
\end{align}

Term (A) is exactly the Wald-RA contribution: the model-predicted contrast at the
individual's covariate value. Term (B) is the correction: it takes the residual
$Y_i - \muYhat(Z_i, X_i)$ --- the part of the outcome that the model missed ---
and reweights it by the inverse instrument propensity score. The AIPW estimator
of the reduced form is then simply
\begin{align}
  \hat{\gamma}_Y^{\,\mathrm{AIPW}}
  = \frac{1}{N}\sum_{i=1}^{N} \YRF_i.
  \label{eq:aipw_estimator}
\end{align}

\paragraph{Algebraic rearrangement.}
It is instructive to rewrite Term~(B) by adding and subtracting $\ehat(X_i)$
in the numerators:
\begin{align}
  \text{(B)}
  &=
  \frac{Z_i - \ehat(X_i)}{\ehat(X_i)}\bigl(Y_i - \muYhat(1,X_i)\bigr)
  +
  \frac{(1-Z_i) - (1-\ehat(X_i))}{1-\ehat(X_i)}\bigl(Y_i - \muYhat(0,X_i)\bigr)
  \cdot(-1)
  \notag \\
  &=
  \frac{Z_i - \ehat(X_i)}{\ehat(X_i)(1-\ehat(X_i))}
  \bigl(Y_i - \muYhat(Z_i, X_i)\bigr).
  \label{eq:B_rewrite}
\end{align}
In this form, Term~(B) is the product of two residuals:
\begin{align}
  \underbrace{Z_i - \ehat(X_i)}_{\text{instrument residual}}
  \quad \times \quad
  \underbrace{Y_i - \muYhat(Z_i,X_i)}_{\text{outcome residual}},
  \label{eq:product_residuals}
\end{align}
scaled by $[\ehat(X_i)(1-\ehat(X_i))]^{-1}$. The double-robustness argument reduces
to showing that this product has zero expectation if either factor has zero conditional
mean.

% ------------------------------------------------------------
\subsection{Proof of Double Robustness: Case 1}
\label{subsec:dr_case1}
% ------------------------------------------------------------

\begin{proposition}[Case~1: Outcome model correct, propensity score possibly wrong]
\label{prop:dr_case1}
Suppose $\muYhat(z,x) = \muY(z,x)$ for all $z \in \{0,1\}$ and all $x$ in the
support of $X_i$, but allow $\ehat(x) \neq e(x)$ arbitrarily. Then
$\hat{\gamma}_Y^{\,\mathrm{AIPW}} \xrightarrow{p} \gamma_Y$.
\end{proposition}

\begin{proof}
Since $\muYhat = \muY$ by assumption, the outcome residual in Term~(B) satisfies
\begin{align}
  Y_i - \muYhat(Z_i, X_i)
  = Y_i - \muY(Z_i, X_i)
  =: \varepsilon_i,
  \label{eq:case1_residual}
\end{align}
where by definition of $\muY$ we have $\E[\varepsilon_i \mid Z_i, X_i] = 0$. Now
compute the expectation of Term~(B) using the law of iterated expectations:
\begin{align}
  \E\!\left[
    \frac{Z_i - \ehat(X_i)}{\ehat(X_i)(1-\ehat(X_i))}\,\varepsilon_i
  \right]
  &= \E\!\left[
      \frac{Z_i - \ehat(X_i)}{\ehat(X_i)(1-\ehat(X_i))}
      \underbrace{\E[\varepsilon_i \mid Z_i, X_i]}_{=\,0}
    \right]
  = 0.
  \label{eq:case1_zero}
\end{align}
The second step uses the fact that $\ehat(X_i)$ is a function of $X_i$ alone, so it
is measurable with respect to the conditioning sigma-algebra. Hence Term~(B) has
zero expectation regardless of $\ehat$, and
\begin{align}
  \E[\YRF_i]
  = \E[\muY(1,X_i) - \muY(0,X_i)] + 0
  = \gamma_Y.
  \label{eq:case1_result}
\end{align}
The sample average $\hat{\gamma}_Y^{\,\mathrm{AIPW}} = \En[\YRF_i]$ therefore
converges in probability to $\gamma_Y$ by the law of large numbers.
\end{proof}

\begin{intuition}{What happens when the outcome model is perfect}
When $\muYhat = \muY$ exactly, the outcome residual $Y_i - \muYhat(Z_i,X_i)$ is
pure noise with mean zero conditional on any function of $(Z_i,X_i)$, including
$\ehat(X_i)$. The IPW correction term is therefore multiplying a mean-zero variable
by a weight, and the product averages to zero regardless of what the weights are.
The Wald-AIPW collapses to the Wald-RA, which is consistent. The propensity score
\emph{drops out entirely}.
\end{intuition}

% ------------------------------------------------------------
\subsection{Proof of Double Robustness: Case 2}
\label{subsec:dr_case2}
% ------------------------------------------------------------

\begin{proposition}[Case~2: Propensity score correct, outcome model possibly wrong]
\label{prop:dr_case2}
Suppose $\ehat(x) = e(x)$ for all $x$ in the support of $X_i$, but allow
$\muYhat(z,x) \neq \muY(z,x)$ arbitrarily. Then
$\hat{\gamma}_Y^{\,\mathrm{AIPW}} \xrightarrow{p} \gamma_Y$.
\end{proposition}

\begin{proof}
We evaluate the full pseudo-outcome $\YRF_i$ in expectation. Start with the
$Z=1$ component of Term~(B):
\begin{align}
  \E\!\left[
    \frac{Z_i\bigl(Y_i - \muYhat(1,X_i)\bigr)}{e(X_i)}
  \right]
  &=
  \E\!\left[
    \frac{1}{e(X_i)}
    \E\!\bigl[Z_i\bigl(Y_i - \muYhat(1,X_i)\bigr) \;\big|\; X_i\bigr]
  \right].
  \label{eq:case2_step1}
\end{align}
Conditioning on $X_i$, we use $\E[Z_i \mid X_i] = e(X_i)$ and---crucially---the
fact that $\muYhat(1,X_i)$ depends only on $X_i$, so it factors out:
\begin{align}
  \E\!\bigl[Z_i\bigl(Y_i - \muYhat(1,X_i)\bigr) \;\big|\; X_i\bigr]
  &= \E[Z_i Y_i \mid X_i] - \muYhat(1,X_i)\,\E[Z_i \mid X_i]
  \notag \\
  &= \E[Z_i Y_i \mid X_i] - \muYhat(1,X_i)\,e(X_i).
  \label{eq:case2_step2}
\end{align}
Next, decompose $\E[Z_i Y_i \mid X_i]$ by conditioning further on $Z_i$:
\begin{align}
  \E[Z_i Y_i \mid X_i]
  &= \Pr(Z_i=1\mid X_i)\,\E[Y_i \mid Z_i=1, X_i]
  + \Pr(Z_i=0\mid X_i)\,\E[Y_i\cdot 0 \mid Z_i=0, X_i]
  \notag \\
  &= e(X_i)\,\muY(1,X_i).
  \label{eq:case2_step3}
\end{align}
Substituting \eqref{eq:case2_step3} into \eqref{eq:case2_step2} and then into
\eqref{eq:case2_step1}:
\begin{align}
  \E\!\left[
    \frac{Z_i\bigl(Y_i - \muYhat(1,X_i)\bigr)}{e(X_i)}
  \right]
  &=
  \E\!\left[
    \frac{e(X_i)\,\muY(1,X_i) - \muYhat(1,X_i)\,e(X_i)}{e(X_i)}
  \right]
  \notag \\
  &=
  \E\!\left[\muY(1,X_i) - \muYhat(1,X_i)\right].
  \label{eq:case2_step4}
\end{align}
The $e(X_i)$ cancels exactly. An identical argument for the $Z=0$ component gives
\begin{align}
  \E\!\left[
    \frac{(1-Z_i)\bigl(Y_i - \muYhat(0,X_i)\bigr)}{1-e(X_i)}
  \right]
  =
  \E\!\left[\muY(0,X_i) - \muYhat(0,X_i)\right].
  \label{eq:case2_step5}
\end{align}
Now assemble the full expected pseudo-outcome by combining Term~(A) with the
two components of Term~(B):
\begin{align}
  \E[\YRF_i]
  &=
  \underbrace{\E\bigl[\muYhat(1,X_i) - \muYhat(0,X_i)\bigr]}_{\text{Term (A)}}
  \notag \\
  &\quad
  +
  \underbrace{
    \E\bigl[\muY(1,X_i) - \muYhat(1,X_i)\bigr]
  }_{\text{correction for $Z=1$ bias}}
  -
  \underbrace{
    \E\bigl[\muY(0,X_i) - \muYhat(0,X_i)\bigr]
  }_{\text{correction for $Z=0$ bias}}
  \notag \\
  &=
  \E\bigl[\muYhat(1,X_i)\bigr]
  + \E\bigl[\muY(1,X_i) - \muYhat(1,X_i)\bigr]
  \notag \\
  &\quad
  -\E\bigl[\muYhat(0,X_i)\bigr]
  - \E\bigl[\muY(0,X_i) - \muYhat(0,X_i)\bigr]
  \notag \\
  &=
  \E\bigl[\muY(1,X_i)\bigr] - \E\bigl[\muY(0,X_i)\bigr]
  = \gamma_Y.
  \label{eq:case2_result}
\end{align}
The model bias $\E[\muYhat(z,X_i) - \muY(z,X_i)]$ is exactly measured and removed by
the IPW correction term, for both $z=0$ and $z=1$. Consistency follows by the law of
large numbers.
\end{proof}

\begin{intuition}{What the correction term is measuring when the propensity score is correct}
Term~(B) is doing something remarkable. It takes the outcome residual
$Y_i - \muYhat(Z_i,X_i)$ --- the part of $Y_i$ that the model missed for the
individual's actual instrument value --- and reweights it to represent what the
model would have missed on average across the whole covariate distribution. Specifically,
$\E[Z_i(Y_i - \muYhat(1,X_i))/e(X_i)]$ equals $\E[\muY(1,X_i) - \muYhat(1,X_i)]$
--- precisely the average bias of the outcome model at $z=1$. Adding this correction
to Term~(A) cancels the model bias term by term. The IPW weights play an essential
role: without them, the residuals $Z_i(Y_i-\muYhat(1,X_i))$ would average only over
the subgroup $\{Z_i=1\}$, which may have a different covariate distribution. The
weights scale each residual up by $1/e(X_i)$ to represent the full population.
\end{intuition}

% ============================================================
\section{The Double Robustness Property: Summary}
\label{sec:dr_summary}
% ============================================================

Propositions~\ref{prop:dr_case1} and~\ref{prop:dr_case2} together establish the
following central result.

\begin{corollary}[Double robustness of the AIPW reduced form]
\label{cor:dr}
Let $\hat{\gamma}_Y^{\,\mathrm{AIPW}} = \En[\YRF_i]$ with $\YRF_i$ defined in
\eqref{eq:rf_pseudo_main}. Then
\begin{align}
  \hat{\gamma}_Y^{\,\mathrm{AIPW}} \;\xrightarrow{p}\; \gamma_Y
\end{align}
if \emph{either} $\muYhat(z,\cdot) \xrightarrow{p} \muY(z,\cdot)$ for $z \in \{0,1\}$,
\emph{or} $\ehat(\cdot) \xrightarrow{p} e(\cdot)$, or both.
\end{corollary}

The same double robustness holds for the first-stage pseudo-outcome $\DFS_i$,
which is defined identically to \eqref{eq:rf_pseudo_main} with $D_i$ replacing $Y_i$
and $\muDhat$ replacing $\muYhat$. Since the Wald-AIPW estimator is the ratio
$\hat{\tau}_{\mathrm{Wald\text{-}AIPW}} = \En[\YRF_i] / \En[\DFS_i]$, and both
numerator and denominator are doubly robust, consistency of the ratio follows whenever
each component is consistent.

\paragraph{A concise view of what cancels.}

Table~\ref{tab:dr_summary} gives a compact overview of the double-robustness argument.

\begin{table}[h!]
\centering
\caption{What cancels in each robustness case}
\label{tab:dr_summary}
\smallskip
\begin{tabular}{lllp{6.5cm}}
\toprule
\textbf{Case} & $\muYhat$ correct? & $\ehat$ correct? & \textbf{Why consistent} \\
\midrule
Wald-RA only   & \checkmark & $\times$
  & Term (B) $\to 0$ because outcome residual
    $\E[\varepsilon_i \mid Z_i,X_i] = 0$ kills it. \\[4pt]
Wald-IPW only  & $\times$ & \checkmark
  & Term (B) exactly measures and subtracts the
    bias of Term (A); the $e(X_i)$ cancel in \eqref{eq:case2_step4}. \\[4pt]
Both correct   & \checkmark & \checkmark
  & Both mechanisms operate; no bias from either source. \\[4pt]
Neither correct & $\times$ & $\times$
  & No cancellation; estimator is generally inconsistent. \\
\bottomrule
\end{tabular}
\end{table}

% ============================================================
\section{Comparison with the Single-Model Alternatives}
\label{sec:comparison}
% ============================================================

\subsection{Why Wald-RA Fails When the Model Is Wrong}
\label{subsec:ra_failure}

Returning to \eqref{eq:ra_bias}: the Wald-RA estimator converges to
$\gamma_Y + \E[b(1,X_i) - b(0,X_i)]$ where $b(z,x) = \muYhat(z,x) - \muY(z,x)$.
This bias is a population-level property of the model misspecification: it is not
reduced by increasing $N$, and it cannot be diagnosed easily because the true
$\muY(z,\cdot)$ is unobserved. In the AIPW estimator, exactly this quantity
$\E[\muYhat(z,X_i) - \muY(z,X_i)]$ is measured and removed by Term~(B) whenever
$\ehat$ is correct (see \eqref{eq:case2_step4}--\eqref{eq:case2_result}).

\subsection{Why Wald-IPW Fails When the Propensity Score Is Wrong}
\label{subsec:ipw_failure}

From \eqref{eq:ipw_bias}, the Wald-IPW estimator converges to
$\E[\muY(1,X_i) \cdot e(X_i)/\ehat(X_i)] - \E[\muY(0,X_i) \cdot (1-e(X_i))/(1-\ehat(X_i))]$.
This differs from $\gamma_Y$ whenever the ratio $e(X_i)/\ehat(X_i)$ is correlated
with $\muY(1,X_i)$ across the covariate distribution, which is generically the case.
In the AIPW estimator, the role of the outcome model is to dampen this sensitivity:
because Term~(A) already captures the systematic variation in $Y_i$ as a function of
$X_i$, the residual $Y_i - \muYhat(Z_i,X_i)$ that enters Term~(B) has smaller
variance than $Y_i$ itself. Even if $\ehat$ is slightly off, the weight distortion
acts on a smaller, mean-zero quantity, limiting the damage.

% ============================================================
\section{Why This Matters for the Thesis}
\label{sec:thesis_relevance}
% ============================================================

The double-robustness argument is not merely a theoretical curiosity. In the empirical
applications considered in this thesis---\citet{angrist1990lifetime},
\citet{card1995using}, \citet{angrist1998children}---the instrument propensity score
$e(X_i) = \Pr(Z_i=1 \mid X_i)$ has a specific structure that makes it comparatively
easy to estimate: the draft lottery number, sibling sex composition, and college
proximity all have near-random variation conditional on $X_i$, so even a flexible
nonparametric model should approximate $e(\cdot)$ well. The outcome equation
$\muY(z,X_i) = \E[Y_i \mid Z_i=z, X_i]$ --- where $Y_i$ is log wages or labour force
participation --- is a much harder object: it reflects the compound influence of family
background, local labour markets, individual ability, and cohort effects, none of which
enter in a simple additive or linear way.

In this asymmetric setting, double robustness is practically valuable: even if the
outcome model is imperfectly estimated, a well-estimated propensity score ensures
that the correction in Term~(B) points in the right direction and removes the
dominant part of the model bias. Conversely, the outcome model provides a low-variance
baseline for Term~(A), reducing the reliance on extreme IPW weights in regions where
$e(X_i)$ is close to zero or one.

Furthermore, the DML implementation of the Wald-AIPW \citep{chernozhukov2018double}
replaces both parametric models with cross-fitted machine learning estimators. This
simultaneously reduces the risk of structural misspecification in either nuisance model
and---by virtue of the Neyman-orthogonality of the AIPW score---preserves valid
$\sqrt{N}$-based inference. The combination of double robustness (protection against
a misspecified model) and Neyman orthogonality (protection against slow ML convergence
rates) is what distinguishes the Wald-AIPW from its single-model alternatives and
motivates its use as the primary DML-based LATE estimator in this thesis.

% ============================================================
% Bibliography
% ============================================================
\bibliographystyle{apalike}
\bibliography{references}

\end{document}
